\begin{trapbox}
	\n
	\begin{description}[style=nextline, leftmargin=2.8em, labelsep=0.5em, itemsep=0.35em]
		\n
		\item[\textbf{\textcolor{CTrap}{易错点一（符号混淆）}}]
			在椭圆中错误使用$c^{2}=a^{2}+b^{2}$，或在双曲线中误用$c^{2}=a^{2}-b^{2}$。\n
		\item[\textbf{\textcolor{CTrap}{正确理解（区分曲线）：}}]
			椭圆关系：$c^{2}=a^{2}-b^{2}$；双曲线关系：$c^{2}=a^{2}+b^{2}$。\n
		\item[\textbf{\textcolor{CTrap}{错因分析（记忆错误）：}}]
			对两种曲线方程推导不熟悉，死记硬背导致符号混乱。\n
	\end{description}
	\n\n

	\medskip
	\n
	\begin{description}[style=nextline, leftmargin=2.8em, labelsep=0.5em, itemsep=0.35em]
		\n
		\item[\textbf{\textcolor{CTrap}{易错点二（滥用公式）}}]
			求抛物线离心率时套用$e=c/a$，但抛物线没有半焦距$c$。\n
		\item[\textbf{\textcolor{CTrap}{正确理解（定义优先）：}}]
			抛物线离心率恒为 $1$，由定义 $e = \frac{|PF|}{d(P,l)}$，且 $|PF| = d(P,l)$，故 $e=1$。\n
		\item[\textbf{\textcolor{CTrap}{错因分析（思维定势）：}}]
			受椭圆、双曲线影响，忽视抛物线定义的独特性。\n
	\end{description}
	\n\n

	\medskip
	\n
	\begin{description}[style=nextline, leftmargin=2.8em, labelsep=0.5em, itemsep=0.35em]
		\n
		\item[\textbf{\textcolor{CTrap}{易错点三（准线公式混淆）}}]
			写准线方程时出现$x=\pm\frac{a}{c}$或$x=\pm\frac{c}{a}$。\n
		\item[\textbf{\textcolor{CTrap}{正确理解（形式正确）：}}]
			准线方程为$x=\pm\frac{a^{2}}{c}$（标准形式下）或$x=\pm\frac{a}{e}$。\n
		\item[\textbf{\textcolor{CTrap}{错因分析（推导缺失）：}}]
			没有理解准线方程来源，仅凭记忆导致错误。\n
	\end{description}
	\n
\end{trapbox}
